%\section{Introduction}
%\ednote{POD to MK: bibliography has some very long entries, lots of space could be gained there}
% \ednote{POD: I think this fits here. NT: wondering if the next three paragraph would not fit better
%   in the odk section; the reader more interested by the MitM part may
%   be put off by this long discussion}

%In mathematics, computer-aided experiments and the use of databases relying on computer calculations such as the Small Groups Library in \GAP~\cite{gap}, the Modular Atlas in group and representation theory, or the $L$-functions and Modular Forms Database (\LMFDB~\cite{lmfdb}), are part of the standard toolbox of the pure mathematician. Certain areas of mathematics completely depend on these libraries. Computers are also increasingly used to support collaborative work and education.

\begin{disclaimer*}
	The following two sections have been previously published as part of \cite{DehKohKon:iop16} with coauthors Paul-Olivier Dehaye, Mihnea Iancu, Michael Kohlhase, Alexander Konovalov, Samuel Leli{\`e}vre, Markus Pfeiffer, Florian Rabe, Nicolas M. Thi{\'e}ry and Tom Wiesing.
	
	Neither the theoretical results nor the majority of the writing can be attributed to me personally. These chapters should hence not be considered my contribution, and is mainly included because it is the best description of the \ODK project and the Math-in-the-Middle approach, which in turn represents a prime concrete application for the results of this thesis.
	
	The Math-in-the-Middle \emph{ontology} for formal mathematics however was largely developed and curated by me, and can thus be considered my contribution.
\end{disclaimer*}

As with interactive theorem provers specifically, in the last decades we witnessed the emergence 
of a wide ecosystem of open-source tools to support research in pure mathematics in general. 
This ranges from specialized to general purpose
computational tools such as \GAP~\cite{gap}, \PariGP, \Linbox, \MPIR, \Sage~\cite{sage}, or \Singular~\cite{singular:on}, via online
databases like the \LMFDB~\cite{lmfdb} or online services like Wikipedia,
\Arxiv~\cite{arxiv}, to webpages like MathOverflow.
A great opportunity is the rapid emergence of key technologies,
in particular the \Jupyter~\cite{jupyter-project:on} (previously \IPython) platform for interactive and
exploratory computing which targets all areas of science.

This has proven the viability and power of collaborative open-source development models,
by users and for users, even for delivering general purpose systems targeting large
audiences such as  researchers, teachers, engineers, amateurs, and others.
Yet some critical long term investments, in particular on the technical side, are in order
to boost the productivity and lower the entry barrier:
\begin{compactitem}
\item Streamlining access, distribution, portability on a wide range of platforms, including
  High Performance Computers or cloud services.
\item Improving user interfaces, in particular in the promising area of collaborative
  workspaces as those provided by CoCalc~\cite{cocalc:on} (previously called \SMC).
\item Lowering barriers between research communities and promoting dissemination. For example
  make it easy for a specialist of scientific computing to use tools from pure
  mathematics, and vice versa.
\item Bringing together the developer communities to promote tighter collaboration and
  symbiosis, accelerate joint development, and share best practices.
\item Structure the development to outsource as much of it as possible to larger communities,
  and focus manpower on core specialities: the implementation of mathematical algorithms
  and databases.
\item And last but not least: Promoting collaborations at all scales to further improve
  the productivity of researchers in pure mathematics and applications.
\end{compactitem}

\ODK\, -- ``Open Digital Research Environment Toolkit for the Advancement
of Mathematics'' \cite{OpenDreamKit:on} -- is a project funded under the
European H2020 Infrastructure call \cite{EINFRA-9} on \emph{Virtual
  Research Environments}, to work on many of these problems.

%In Section~\ref{sec:odk}, we will introduce the \ODK project  to establish the context for the
%``Math-in-the-Middle'' (MitM) integration approach described in
%Section~\ref{sec:mitm}. The remaining sections then elucidate the approach by presenting
%first experiments and refinements of the chosen integration paradigm:
%Section~\ref{sec:lmfdb} details how existing knowledge representation and data structures
%can be represented as MitM interface theories with a case study of equipping the \LMFDB
%with a MitM-based programming interface.  Section~\ref{sec:gapsage} discusses system
%integration between \GAP and \Sage and how this can be routed through a MitM
%ontology. Section~\ref{sec:concl} concludes the paper and discusses future work.

% \ednote{R1: Throughout the paper is would be better to reduce the use of acronyms. 
% For example, the paper is unnecessarily difficult to read because of the extensive 
% use of "VRE" and "COMs" as well as the many other acronyms. AK: I suggest to use VRE,
% but revise the others. POD: I have removed COM and many others, leaving this as a reminder to be on the lookout.}

%%% Local Variables:
%%% mode: latex
%%% TeX-master: "paper"
%%% End:

%  LocalWords:  specialized Arxiv Jupyter IPython ldots compactitem emph mitm lfmdb concl
%  LocalWords:  gaptypes
